The game uses a finite connected undirected graph. C chooses a starting dot, R chooses after seeing C, and C moves first. Turns alternate; each player may stay or cross one road. Positions are visible. Sharing a dot is immediate capture, including at placement. An evasion theorem means a response strategy forever, not survival of a classroom timer.

\textbf{Every finite tree is cop-win.} C moves toward R along their unique path. Delete C's current dot and look at R's component. Each uncaptured cop advance makes that component strictly smaller, and R cannot cross C to escape it. Finiteness gives capture. On an $n$-dot path the best worst-case start is a middle dot, with $\lfloor n/2\rfloor$ cop moves: a farthest endpoint robber who stays gives the lower bound, and advancing from the middle gives the upper bound.

\textbf{A bare cycle with at least four dots is robber-win.} Before each cop turn R can remain outside C's closed neighborhood (C and its neighbors). After C moves, R stays if safe, or moves away to the other ring neighbor, which is safe. A triangle is different: one cop is already adjacent to every dot. The existence of a cycle inside a larger graph alone does not decide who wins.

\textbf{A dominated dot may be folded away.} If $N[v]\subseteq N[u]$ for distinct dots, deleting $v$ preserves cop-win status; staying counts in these closed neighborhoods. A sequence of folds to one dot certifies a cop strategy. The full finite classification says a connected graph is cop-win exactly when it can be reduced this way. This deeper result is an adult destination, not a prerequisite for initial games. The explanation and limits appear with the fold solution.

One diagonal makes a square cop-win; a bare pentagon needs two added chords. One pentagon chord still leaves a square as a retract. Two cops catch on every finite cycle: one blocks a dot while the other sweeps the remaining path. These statements do not claim two cops suffice on every graph.
